\def\tituloconcepto{Algoritmo HHL}
% Si el capitulo se compila solo (sin main.tex), esta cabecera abre el documento.
% Antes de % Si el capitulo se compila solo (sin main.tex), esta cabecera abre el documento.
% Antes de % Si el capitulo se compila solo (sin main.tex), esta cabecera abre el documento.
% Antes de \input{cabecera} cada capitulo define \tituloconcepto.
\ifdefined\enmain\else
  \documentclass[11pt,a4paper]{article}
  \input{preambulo}
  \begin{document}
  \title{\tituloconcepto}
  \author{\autor}
  \date{\today}
  \maketitle
\fi
 cada capitulo define \tituloconcepto.
\ifdefined\enmain\else
  \documentclass[11pt,a4paper]{article}
  % Preambulo compartido por todos los documentos. Compilar con XeLaTeX/LuaLaTeX
% (o Tectonic). Si se usa Overleaf, elegir el compilador XeLaTeX.
\usepackage[spanish,es-noshorthands]{babel}
\usepackage{amsmath,amssymb}
\usepackage{booktabs,array}
\usepackage[a4paper,margin=2.5cm]{geometry}
\usepackage{xcolor}
\usepackage{listings}
\usepackage[most]{tcolorbox}
\usepackage{quantikz}
\usepackage[hidelinks]{hyperref}

% ---- Autor (editar a gusto) -------------------------------------------------
\newcommand{\autor}{Genesis Silesky}

% ---- Notacion de Dirac y atajos ---------------------------------------------
% quantikz ya define \ket y \bra; DeclareRobustCommand los redefine sin error.
\DeclareRobustCommand{\ket}[1]{\left|#1\right\rangle}
\DeclareRobustCommand{\bra}[1]{\left\langle#1\right|}
\DeclareRobustCommand{\braket}[2]{\left\langle#1\middle|#2\right\rangle}
\newcommand{\C}{\mathbb{C}}
\newcommand{\R}{\mathbb{R}}
\DeclareMathOperator{\Tr}{Tr}

% ---- Cajas de apoyo --------------------------------------------------------
% "concepto": definiciones previas que hacen falta para entender el capitulo.
\newtcolorbox{concepto}[1]{breakable,colback=blue!4,colframe=blue!55!black,
  fonttitle=\bfseries,title={Concepto previo: #1}}
% "nota": aclaraciones y advertencias.
\newtcolorbox{nota}[1]{breakable,colback=orange!6,colframe=orange!70!black,
  fonttitle=\bfseries,title={#1}}

% ---- Estilos de listados ----------------------------------------------------
\lstdefinestyle{py}{language=Python,basicstyle=\ttfamily\footnotesize,
  breaklines=true,postbreak=\mbox{\textcolor{gray}{$\hookrightarrow$}\space},
  numbers=left,numberstyle=\tiny\color{gray},frame=single,
  commentstyle=\color{green!40!black}\itshape,
  keywordstyle=\color{blue!70!black}\bfseries,
  stringstyle=\color{red!60!black},showstringspaces=false,tabsize=4,
  columns=fullflexible,keepspaces=true}
\lstdefinestyle{salida}{basicstyle=\ttfamily\footnotesize,breaklines=true,
  frame=leftline,backgroundcolor=\color{gray!8},columns=fullflexible,
  keepspaces=true}

  \begin{document}
  \title{\tituloconcepto}
  \author{\autor}
  \date{\today}
  \maketitle
\fi
 cada capitulo define \tituloconcepto.
\ifdefined\enmain\else
  \documentclass[11pt,a4paper]{article}
  % Preambulo compartido por todos los documentos. Compilar con XeLaTeX/LuaLaTeX
% (o Tectonic). Si se usa Overleaf, elegir el compilador XeLaTeX.
\usepackage[spanish,es-noshorthands]{babel}
\usepackage{amsmath,amssymb}
\usepackage{booktabs,array}
\usepackage[a4paper,margin=2.5cm]{geometry}
\usepackage{xcolor}
\usepackage{listings}
\usepackage[most]{tcolorbox}
\usepackage{quantikz}
\usepackage[hidelinks]{hyperref}

% ---- Autor (editar a gusto) -------------------------------------------------
\newcommand{\autor}{Genesis Silesky}

% ---- Notacion de Dirac y atajos ---------------------------------------------
% quantikz ya define \ket y \bra; DeclareRobustCommand los redefine sin error.
\DeclareRobustCommand{\ket}[1]{\left|#1\right\rangle}
\DeclareRobustCommand{\bra}[1]{\left\langle#1\right|}
\DeclareRobustCommand{\braket}[2]{\left\langle#1\middle|#2\right\rangle}
\newcommand{\C}{\mathbb{C}}
\newcommand{\R}{\mathbb{R}}
\DeclareMathOperator{\Tr}{Tr}

% ---- Cajas de apoyo --------------------------------------------------------
% "concepto": definiciones previas que hacen falta para entender el capitulo.
\newtcolorbox{concepto}[1]{breakable,colback=blue!4,colframe=blue!55!black,
  fonttitle=\bfseries,title={Concepto previo: #1}}
% "nota": aclaraciones y advertencias.
\newtcolorbox{nota}[1]{breakable,colback=orange!6,colframe=orange!70!black,
  fonttitle=\bfseries,title={#1}}

% ---- Estilos de listados ----------------------------------------------------
\lstdefinestyle{py}{language=Python,basicstyle=\ttfamily\footnotesize,
  breaklines=true,postbreak=\mbox{\textcolor{gray}{$\hookrightarrow$}\space},
  numbers=left,numberstyle=\tiny\color{gray},frame=single,
  commentstyle=\color{green!40!black}\itshape,
  keywordstyle=\color{blue!70!black}\bfseries,
  stringstyle=\color{red!60!black},showstringspaces=false,tabsize=4,
  columns=fullflexible,keepspaces=true}
\lstdefinestyle{salida}{basicstyle=\ttfamily\footnotesize,breaklines=true,
  frame=leftline,backgroundcolor=\color{gray!8},columns=fullflexible,
  keepspaces=true}

  \begin{document}
  \title{\tituloconcepto}
  \author{\autor}
  \date{\today}
  \maketitle
\fi


\section{Algoritmo HHL: resolver un sistema lineal con dos qubits}

\subsection{Objetivo}
Entender cómo el algoritmo HHL resuelve $A\mathbf x=\mathbf b$ dividiendo los
coeficientes de $\mathbf b$ entre los autovalores de $A$, y construir su versión
de dos qubits para una matriz $2\times2$.

\subsection{Conceptos previos}
\begin{concepto}{autovalores, autovectores y estimación de fase}
Ver capítulo~17 (estimación de fase) y la sección~\ref{sec:prerreq}. Una matriz
hermítica $A$ se escribe como $A=\sum_i\lambda_i\ket{u_i}\bra{u_i}$, con
autovalores $\lambda_i$ y autovectores $\ket{u_i}$.
\end{concepto}

\subsection{Fundamento teórico}
HHL (Harrow, Hassidim y Lloyd, 2009) resuelve sistemas de ecuaciones lineales
$A\mathbf x=\mathbf b$, con $A$ una matriz hermítica $N\times N$. Clásicamente, con
métodos directos o iterativos, el tiempo es $O(N^3)$; en condiciones ideales,
HHL lo reduce a $O(\log N)$. Se usa como subrutina en modelos de aprendizaje
automático cuántico.

\paragraph{Un ejemplo clásico.} Dos grupos hacen vasijas. En 2.5 horas, el grupo
A hace una vasija más que el grupo B en 0.5 horas; y en 0.5 horas, A hace tantas
como B en 2.5 horas. Si $x$ e $y$ son las vasijas por hora:
\begin{equation}
  \begin{pmatrix}2.5&-0.5\\-0.5&2.5\end{pmatrix}
  \binom xy=\binom10 ,\qquad x=\tfrac5{12},\ y=\tfrac1{12}.
\end{equation}

\paragraph{La idea de HHL.} Los vectores se codifican como estados cuánticos
normalizados, $\ket b$ y $\ket x$, y resolver $\mathbf x=A^{-1}\mathbf b$ pasa a ser
$\ket x\propto A^{-1}\ket b$. Se escribe $\ket b$ en la base de autovectores,
$\ket b=\sum_i\beta_i\ket{u_i}$ con $\beta_i=\langle u_i|b\rangle$, y entonces
\begin{equation}
  \ket x\ \propto\ \sum_i\frac{\beta_i}{\lambda_i}\ket{u_i}.
\end{equation}
Basta dividir cada coeficiente entre su autovalor. Para la matriz del
ejemplo, los autovalores son $\lambda_1=2$ y $\lambda_2=3$, con autovectores
$\ket{u_1}=\frac1{\sqrt2}\binom11$ y $\ket{u_2}=\frac1{\sqrt2}\binom1{-1}$.

\paragraph{Los tres componentes.} HHL usa tres grupos de qubits: de codificación
(guardan $\ket b$, $n_b=\log_2N$ qubits), de registro (guardan los autovalores
estimados) y un qubit auxiliar. Los pasos son:
\begin{enumerate}
  \item \textbf{Estimación de fase:} estima los autovalores de $A$ y los guarda en el
    registro: $\sum_i\beta_i\ket{\tilde\lambda_i}\ket{u_i}$.
  \item \textbf{Rotación controlada:} usando el registro como control, gira el
    qubit auxiliar un ángulo $\theta_i=2\sin^{-1}(c/\lambda_i)$. Así el
    recíproco del autovalor pasa a la amplitud del auxiliar:
    $\sum_i\beta_i\ket{\lambda_i}\ket{u_i}\big(\sqrt{1-c^2/\lambda_i^2}\ket0
    +\frac c{\lambda_i}\ket1\big)$. El parámetro $c$ cumple $c\le\min|\lambda_i|$.
  \item \textbf{Estimación de fase inversa:} desentrelaza el registro.
  \item \textbf{Medición:} si el auxiliar vale $\ket1$ (con cierta probabilidad),
    el qubit de codificación queda en $\ket x\propto\sum_i\frac{\beta_i}{\lambda_i}\ket{u_i}$.
\end{enumerate}

\paragraph{Versión de dos qubits.} Con una matriz $2\times2$ de autovalores
$2$ y $3$, el registro necesita un solo qubit y el circuito se simplifica
hasta quedar en dos qubits: un qubit de codificación y uno auxiliar. Si
$A=U\,\mathrm{diag}(\lambda_1,\lambda_2)\,U^\dagger$, el circuito es:
$U^\dagger$ sobre el qubit de codificación (los autovectores pasan a $\ket0$ y
$\ket1$), una rotación $R_y(\theta_i)$ del auxiliar controlada por el qubit de
codificación y, al final, $U$. Con $\lambda_1=2$ y $\lambda_2=3$, el mayor valor
posible de $c$ es $2$.

\subsection{Código (completamente comentado)}
\lstinputlisting[style=py]{../../fuentes/18_hhl.py}

\subsection{Explicación celda por celda}
\begin{description}
  \item[Celda 2] Resuelve el sistema de las vasijas de forma clásica.
  \item[Celda 3] Calcula autovalores y autovectores, los coeficientes
    $\beta_i$, y comprueba que $\sum_i(\beta_i/\lambda_i)\ket{u_i}$ da la misma
    solución que la clásica (normalizada).
  \item[Celda 4] Define el circuito HHL de dos qubits: codificar $\ket b$,
    $U^\dagger$, rotaciones controladas y $U$.
  \item[Celda 5] Calcula las amplitudes del qubit de codificación cuando el
    auxiliar vale $\ket1$ y las compara con la solución clásica.
  \item[Celda 6] Repite con la matriz del experimento y comprueba que su $U$ es
    $R_y(\pi/4)$.
  \item[Celda 7] Imprime las probabilidades de las cuatro cadenas
    $\ket{q_1q_0}$ que se medirían.
\end{description}

\subsection{Salida obtenida}
\lstinputlisting[style=salida]{../../fuentes/salidas/18_hhl.txt}
(El dibujo del circuito no se reproduce aquí.) Es determinista.

\subsection{Interpretación}
Para el sistema de las vasijas, la solución cuántica normalizada es
$(0.9806,0.1961)$, igual que la clásica $(5/12,1/12)$ normalizada. Para la
matriz del experimento, $A=\begin{psmallmatrix}2.14645&-0.35355\\-0.35355&2.85355\end{psmallmatrix}$
y $\mathbf b=(0.70711,0.70711)$, también coinciden: $(0.7887,0.6148)$. Sus
autovalores son $2$ y $3$ y su matriz $U$ coincide con $R_y(\pi/4)$ (salvo
signos). Al medir, la solución solo se obtiene cuando el auxiliar vale $1$,
con probabilidad $0.72$ en el primer caso y $0.92$ en el segundo.

\begin{nota}{Qué se simplifica en este ejemplo}
Igual que en el PDF de la práctica, aquí los autovalores de $A$ se
\emph{conocen} y se usan directamente para armar las rotaciones; la estimación de
fase que los obtiene en un HHL completo se sustituye por esa información previa.
Además, un estado cuántico no se puede leer entero: el código calcula las
amplitudes con el simulador, lo que un equipo real no permite.
\end{nota}

\subsection{Ejercicios}
\begin{enumerate}
  \item Cambia $\mathbf b$ por $(\cos\varphi,\sin\varphi)$ con otro ángulo y compara
    con la solución clásica.
  \item Cambia $c$ a un valor menor que $2$: ¿qué pasa con la probabilidad de que
    el auxiliar valga $1$?
  \item Prueba con una matriz $2\times2$ simétrica distinta (con autovalores
    positivos) y comprueba el resultado.
\end{enumerate}

\subsection{Errores comunes}
\begin{itemize}
  \item Confundir $\ket x$ con $\mathbf x$: HHL entrega el vector
    \emph{normalizado}, no su tamaño.
  \item Elegir $c$ mayor que el menor autovalor: entonces $c/\lambda_i>1$ y el
    ángulo $2\sin^{-1}(c/\lambda_i)$ no existe.
\end{itemize}

\ifdefined\enmain\else
  \end{document}
\fi

